\subsection{Quotient spaces and subspaces of a weak space}

Let F, G be two real vector spaces in duality. Let M be a vector subspace of F, and consider the subspace N of the orthogonal \(M^{\circ}\) in G; if \(y_{1}, y_{2}\) are two points of G that are congruent mod. N then \(\langle x, y_{1} \rangle = \langle x, y_{2} \rangle\) for all \(x \in M\) . For each class \(\dot{y}\) mod. N, denote the common value of \(\langle x, y \rangle\) when y varies in \(\dot{y}\) by \(\langle x, \dot{y} \rangle\) ; clearly \((x, \dot{y}) \mapsto \langle x, \dot{y} \rangle\) is a bilinear form on \(M \times (G/N)\) .

PROPOSITION 7. --- Let M be a vector subspace of F and N a vector subspace of G where F and G are two vector spaces in duality. Suppose that M and N are orthogonal (which is equivalent to saying that \(N \subset M^{\circ}\) , or \(M \subset N^{\circ}\) ). The vector spaces M and G/N are then in duality by the bilinear form \((x, \dot{y}) \mapsto \langle x, \dot{y} \rangle\) .

\begin{enumerate}
\def\labelenumi{(\roman{enumi})}
\item
  The topology \(\sigma(\mathbf{M},\mathbf{G} / \mathbf{N})\) for this duality is induced by \(\sigma(\mathbf{F},\mathbf{G})\) (and in particular we have \(\sigma(\mathbf{F},\mathbf{G}) = \sigma(\mathbf{F},\mathbf{G} / \mathbf{F}^{\circ})\) ).
\item
  The topology \(\sigma(\mathbf{G}/\mathbf{N}, \mathbf{M})\) for this duality is coarser than the quotient topology of \(\sigma(\mathbf{G}, \mathbf{F})\) by \(\mathbf{N}\) ; these topologies are identical if and only if \(\mathbf{M} + \mathbf{G}^{\circ} = \mathbf{N}^{\circ}\) .
\setcounter{enumi}{0}
\item
  Every element of \(\mathbf{G} / \mathbf{N}\) is a class mod. \(\mathbf{N}\) of an element of \(\mathbf{G}\) ; if \(z_{i}\) ( \(1 \leqslant i \leqslant n\) ) are elements of \(\mathbf{G}\) and \(\dot{z}_{i}\) ( \(1 \leqslant i \leqslant n\) ) is the class of \(z_{i}\) in \(\mathbf{G} / \mathbf{N}\) then the set of \(y \in \mathbf{M}\) such that \(|\langle y, \dot{z}_{i} \rangle| \leqslant \alpha\) for \(1 \leqslant i \leqslant n\) is the trace on \(\mathbf{M}\) of the set of those \(x \in \mathbf{F}\) such that \(|\langle x, z_{i} \rangle| \leqslant \alpha\) for \(1 \leqslant i \leqslant n\) ; the conclusion follows from the definition of neighbourhoods of 0 for the weak topology.
\item
  Let \(p: \mathbf{G} \to \mathbf{G} / \mathbf{N}\) be the canonical surjection. We show that the quotient topology \(\mathcal{T}\) of \(\sigma(\mathbf{G}, \mathbf{F})\) by \(\mathbf{N}\) is identical with \(\sigma(\mathbf{G} / \mathbf{N}, \mathbf{N}^{\circ})\) . As, for \(z \in \mathbf{G}\) , \(y \in \mathbf{N}^{\circ}\) , we have \(\langle y, p(z) \rangle = \langle y, z \rangle\) , it follows that every neighbourhood of 0 for \(\sigma(\mathbf{G} / \mathbf{N}, \mathbf{N}^{\circ})\) is of the form \(p(\mathbf{V})\) , where \(\mathbf{V}\) is a neighbourhood of 0 for \(\sigma(\mathbf{G}, \mathbf{F})\) saturated for the relation \(z - z' \in \mathbf{N}\) , therefore \(\mathcal{T}\) is finer than \(\sigma(\mathbf{G} / \mathbf{N}, \mathbf{N}^{\circ})\) . Conversely let \(\mathbf{U} = \mathbf{W}(y_1, ..., y_n; \alpha)\) be a neighbourhood of 0 in \(\mathbf{G}\) for \(\sigma(\mathbf{G}, \mathbf{F})\) , where \(y_i \in \mathbf{F}\) for \(1 \leqslant i \leqslant n\) and \(\alpha > 0\) ; we are going to see that for \(1 \leqslant i \leqslant n\) , there exist elements \(t_i \in \mathbf{N}^{\circ}\) such that if one puts \(\mathbf{U}' = \mathbf{W}(t_1, ..., t_n; \alpha)\) , then \(p(\mathbf{U}') \subset p(\mathbf{U})\) ; this will show that \(\sigma(\mathbf{G} / \mathbf{N}, \mathbf{N}^{\circ})\) is finer than \(\mathcal{T}\) and therefore is actually identical with \(\mathcal{T}\) . Now, let \(L\) be the vector subspace of \(F\) generated by \(N^{\circ}\) and the \(y_i\) , and denote by \(P\) the complementary subspace of \(N^{\circ}\) in \(L\) ; it is of finite dimension, say \(m\) . Let \((x_j)_{1 \leqslant j \leqslant m}\) be a basis of \(P\) ; the restrictions to \(N\) of the linear forms \(x \mapsto \langle x_j, z \rangle\) are linearly independent, since otherwise there exists \(x \neq 0\) in \(P\) such that \(\langle x, z \rangle = 0\) for all \(z \in N\) , that is to say \(x \in N^{\circ}\) , which contradicts the definition of \(P\) . Thus we conclude that for all \(z' \in G\) , there exists \(s \in N\) such that \(\langle x_j, z' \rangle = \langle x_j, s \rangle\) for all \(j\) ; if \(z' = z + s\) , we have \(\langle x, z \rangle = 0\) for all \(x \in P\) . This being so, put \(y_i = t_i + w_i\) , where \(t_i \in N^{\circ}\) and \(w_i \in P\) ; we have \(\langle y_i, z \rangle = \langle t_i, z \rangle = \langle t_i, z' \rangle\) for \(1 \leqslant i \leqslant n\) ; therefore, for all \(z' \in U'\) , there exists \(z \in U\) such that \(z' - z \in N\) , that is to say we have \(p(\mathrm{U}') \subset p(\mathrm{U})\) .
\end{enumerate}

Returning to the case where M is any subspace of \(N^{\circ}\) , note that evidently \(\sigma(\mathrm{G}/\mathrm{N},\mathrm{M})=\sigma(\mathrm{G}/\mathrm{N},\mathrm{M}+\mathrm{G}^{\circ})\) ; further, from prop. 3 of II, p.~43, we see that, if \(y\in N^{\circ}\) is such that the linear form \(\dot{z}\mapsto\langle y,\dot{z}\rangle\) is continuous for \(\sigma(\mathrm{G}/\mathrm{N},\mathrm{M})\) , then necessarily \(y\in M+G^{\circ}\) . We conclude that the condition \(M+G^{\circ}=N^{\circ}\) is necessary and sufficient for the quotient topology T to be equal to \(\sigma(\mathrm{G}/\mathrm{N},\mathrm{M})\) .

Remark. --- The duality between M and G/N (where M and N are two orthogonal subspaces) is separating in M, if and only if \(M \cap G^{\circ} = \{0\}\) ; it is separating in G/N, if and only if \(N = M^{\circ}\) .

COROLLARY 1. --- Suppose that the duality between F and G is separating in F. For a vector subspace M of F the topology \(\sigma(\mathbf{G}/\mathbf{M}^{\circ}, \mathbf{M})\) is identical with quotient topology of \(\sigma(\mathbf{G}, \mathbf{F})\) by \(\mathbf{M}^{\circ}\) , if and only if M is closed for the topology \(\sigma(\mathbf{F}, \mathbf{G})\) .

This follows from prop. 7 putting \(\mathbf{N} = \mathbf{M}^{\circ}\) , and recalling that \(\mathbf{M}^{\circ \circ}\) is the closure of \(\mathbf{M}\) for \(\sigma (\mathrm{F},\mathrm{G})\) (II, p.~45, cor. 2).

COROLLARY 2. --- If M is of finite dimension n and the duality is separating in F, then \(M^{\circ}\) is of codimension n in G. If M is closed for \(\sigma(F, G)\) and of finite codimension n and if the duality is separating in G, then \(M^{\circ}\) is of dimension n.

For, \(G/M^{\circ}\) is in separating duality with M; if M is of dimension n, the same is therefore true of \(G/M^{\circ}\) (II, p.~41, example 1). If M is closed, \(F/M = F/M^{\circ\circ}\) is in separating duality with \(M^{\circ}\) ; if F/M is of dimension n, it is therefore the same for \(M^{\circ}\) (II, p.~41, example 1).

COROLLARY 3. --- Let (F, G), (F₁, G₁) be two pairs of spaces in separating duality and let u be a linear mapping of F in F₁, which is continuous for σ(F, G) and σ(F₁, G₁). Then u is a strict morphism of F in F₁, if and only if, Im(\({}^t u\)) is a closed subspace in G for σ(G, F).

Let \(\mathbf{N} = \operatorname{Im}(^t u) \subset \mathbf{G}\) ; we know that \(\mathrm{N}^{\circ} = \mathrm{Ker}(u)\) in \(\mathbf{F}\) (II, p.~47, formula (3)). Let \(p: \mathbf{F} \to \mathbf{F}/\mathbf{N}^{\circ}\) be the canonical mapping so that \(u\) factorises as

\[
u: \mathrm{F} \xrightarrow {p} \mathrm{F} / \mathrm{N} ^ {\circ} \xrightarrow {w} \mathrm{F} _ {1},
\]

where w is injective. The spaces \(F/N^{\circ}\) and N are in separating duality and by formula (1) of II, p.~48, we have \(\langle w(\dot{y}), z_{1} \rangle = \langle \dot{y}, {}^{t}u(z_{1}) \rangle\) for all \(\dot{y} \in F/N^{\circ}\) and \(z_{1} \in G_{1}\) . This relation shows that w is an isomorphism of \(F/N^{\circ}\) , carrying the topology \(\sigma(F/N^{\circ}, N)\) , on \(u(F)\) with the topology induced by \(\sigma(F_{1}, G_{1})\) . The conclusion results therefore from cor. 1 and the definition of a strict morphism.

COROLLARY 4. --- Let \((\mathrm{F}, \mathrm{G})\) , \((\mathrm{F}_{1}, \mathrm{G}_{1})\) be two pairs in separating duality, and let u be a linear mapping of F in \(F_{1}\) that is continuous for \(\sigma(\mathrm{F}, \mathrm{G})\) and \(\sigma(\mathrm{F}_{1}, \mathrm{G}_{1})\) . Then u is surjective, if and only if, \(^{t}u\) is an isomorphism of \(G_{1}\) (with topology \(\sigma(G_{1}, F_{1})\) ) on \(^{t}u(G_{1})\) with the topology induced by \(\sigma(\mathrm{G}, \mathrm{F})\) .

For, to say that \(u(\mathrm{F}) = \mathrm{F}_{1}\) is equivalent to saying that \(u(\mathrm{F})\) is closed and everywhere dense in \(F_{1}\) for \(\sigma(\mathrm{F}_{1}, \mathrm{G}_{1})\) ; cor. 4 follows then from cor. 3 applied to \(^{t}u\) and of II, p.~47, cor. 2.

Remarks. --- 1) Let \((\mathrm{F}_1, \mathrm{G}_1), (\mathrm{F}_2, \mathrm{G}_2), (\mathrm{F}_3, \mathrm{G}_3)\) be three pairs of spaces in separating duality and consider a sequence of two linear mappings

\[
\mathrm{F} _ {1} \xrightarrow {u} \mathrm{F} _ {2} \xrightarrow {v} \mathrm{F} _ {3}\tag{5}
\]

that are continuous for the weak topologies corresponding respectively with \(\mathbf{G}_1\) , \(\mathbf{G}_2\) , \(\mathbf{G}_3\) ; we consider the sequence of transposed mappings

\[
\mathrm{G} _ {3} \stackrel {{t _ {v}}} {{\to}} \mathrm{G} _ {2} \stackrel {{t _ {u}}} {{\to}} \mathrm{G} _ {1}.\tag{6}
\]

It is clear that \({}^t (v\circ u) = {}^t u\circ {}^tv\) , therefore the relation \(v\circ u = 0\) is equivalent to \({}^t u\circ {}^tv = 0\) . The sequence (5) is exact if, and only if, the three following conditions are satisfied a) \({}^t u\circ {}^tv = 0\) :

\begin{enumerate}
\def\labelenumi{\alph{enumi})}
\setcounter{enumi}{1}
\item
  \(\operatorname{Im}(^{t}v)\) is dense in \(\operatorname{Ker}(^{t}u)\) ;
\item
  \(^t u\) is a strict morphism of \(G_2\) in \(G_1\) .
\end{enumerate}

This follows in effect from cor. 3 of II, p.~49 and formulae (3) and (4) of II, p.~47. 2) It must not be thought that when \(u\) is a strict morphism of \(F\) in \(F_1\) , then \(^t u\) is necessarily a strict morphism of \(G_1\) in \(G\) ; in other words \(u\) can be a strict morphism without \(u(F)\) being closed in \(F_1\) for \(\sigma(F_1, G_1)\) . This is shown by the example where \(F\) is a non-closed subspace of \(F_1\) and \(G = G_1 / F^\circ\) , \(u\) being the canonical injection. Similarly, the fact that the sequence (5) is exact does not necessarily imply that (6) is exact, however, if the sequence (5) is exact and if \(v\) is a strict morphism, then the sequence (6) is exact, by the remark 1 and by II, p.~49, cor. 3.

